% ============================================================
%  GUÍA — INTEGRALES ALGEBRAICAS Y CAMBIO DE VARIABLE
%  Cálculo Integral · Método de sustitución y división de polinomios
%  Fuente de la clase: tex_files/clases_integrales_ruben_usm/memoria_clases.md
%  Autor: Ing. Gabriel Astudillo
%  Compilar:  latexmk -xelatex guia_integrales_algebraicas_cdv.tex
%
%  Alcance: integrales indefinidas algebraicas, cambio de variable y
%  división de polinomios (grado del numerador ≥ grado del denominador).
%  NO incluye: fracciones parciales, integración por partes ni integrales
%  trigonométricas (se mencionan solo como repertorio futuro).
% ============================================================

\documentclass[12pt, letterpaper]{article}

% ── Paquetes ─────────────────────────────────────────────────
\usepackage{mathwizards-guia}
\mwguia{Guía — Integrales Algebraicas y Cambio de Variable}{Ing. Gabriel Astudillo $\cdot$ Cálculo Integral}

% ── Comandos propios de esta guía ────────────────────────────
\newcommand{\dx}{\,dx}
\newcommand{\du}{\,du}
\newcommand{\CC}{\,+\,C}
\newcommand{\R}{\mathbb{R}}

\begin{document}
% ============================================================

% ── Portada / Título ─────────────────────────────────────────
\mwportada{Integrales Algebraicas y Cambio de Variable}{Método de sustitución y división de polinomios}{Guía de estudio $\cdot$ 30 problemas progresivos con respuestas}

\vspace{4pt}
\begin{cuadroinfo}
\small
\textbf{Presentación:} Esta guía trabaja la integración indefinida de expresiones
\emph{algebraicas} por dos caminos que ya conoces de clase: el \textbf{cambio de
variable} y la \textbf{división previa de polinomios}. Los 30 problemas están
ordenados en tres bloques de dificultad creciente, desde sustituciones directas
hasta integrales con radicales anidados y fracciones racionales que exigen
dividir antes de integrar.
\\[4pt]
\textbf{Alcance:} potencias, sumas y restas, cambio de variable y división de
polinomios. \textbf{No} se usan fracciones parciales, integración por partes ni
integrales trigonométricas.
\\[4pt]
\textbf{Leyenda de dificultad:}\quad
\medio{} Intermedio \quad
\dificil{} Desafiante \quad
\muyDificil{} Avanzado
\end{cuadroinfo}

\vspace{8pt}

% ============================================================
\section{Formulario de referencia}
% ============================================================

\subsection{Reglas básicas de integración}

\begin{cuadroteoria}
\begin{align*}
  \int x^{n}\dx &= \frac{x^{n+1}}{n+1} \CC \qquad (n \neq -1) &
  \int \frac{1}{x}\dx &= \ln|x| \CC \\[6pt]
  \int k\,f(x)\dx &= k\int f(x)\dx &
  \int \bigl[f(x) \pm g(x)\bigr]\dx &= \int f(x)\dx \pm \int g(x)\dx \\[6pt]
  \int k\dx &= kx \CC & \int \du &= u \CC
\end{align*}
\end{cuadroteoria}

\subsection{Propiedades de exponentes}

\begin{cuadroteoria}
\textbf{Potencias y radicales} (para pasar de raíz a exponente fraccionario y viceversa):
\[
\sqrt[m]{x^{n}} = x^{n/m}
\qquad
x^{n/m} = \sqrt[m]{x^{n}}
\qquad
a^{m}a^{n} = a^{m+n}
\qquad
\frac{a^{m}}{a^{n}} = a^{m-n}
\]
\textbf{Exponentes negativos y distribución de un denominador} (las usaremos en la §4):
\[
\left(\frac{a}{b}\right)^{-n} = \left(\frac{b}{a}\right)^{n}
\qquad
\frac{1}{a^{n}} = a^{-n}
\qquad
a^{-n} = \frac{1}{a^{n}}
\qquad
\frac{a \pm b \pm c \pm d}{n} = \frac{a}{n} \pm \frac{b}{n} \pm \frac{c}{n} \pm \frac{d}{n}
\]
\end{cuadroteoria}

\subsection{Método de cambio de variable (sustitución)}

\begin{cuadroteoria}
El cambio de variable \emph{deshace} la regla de la cadena. Si $u = g(x)$, entonces
$du = g'(x)\dx$ y
\[
\int f\bigl(g(x)\bigr)\,g'(x)\dx = \int f(u)\du = F(u) \CC,
\qquad F' = f
\]
La integral se reduce a una forma elemental sobre la variable auxiliar $u$.
\end{cuadroteoria}

\begin{cuadropasos}
\textbf{Procedimiento en 5 pasos:}
\begin{enumerate}
  \item \textbf{Elegir} la sustitución $u = g(x)$ (el argumento interno, el denominador
        o la base de una potencia).
  \item \textbf{Diferenciar}: $du = g'(x)\dx$; despejar $\dx$ si conviene.
  \item \textbf{Reescribir} toda la integral en $u$. No debe quedar ninguna $x$; si
        queda, despeja $x$ de $u = g(x)$ y sustitúyela.
  \item \textbf{Integrar} respecto a $u$.
  \item \textbf{Regresar} a $x$ reemplazando $u = g(x)$ y añadir $+C$.
\end{enumerate}
\end{cuadropasos}

\begin{cuadroadvertencia}
\textbf{Errores clásicos:}
\begin{itemize}
  \item Olvidar reemplazar el $\dx$ por su equivalente en $u$ (dejar una $x$ suelta).
  \item No despejar correctamente $x$ cuando la sustitución deja términos sobrantes.
  \item Terminar en $u$ y \emph{no} volver a la variable original $x$.
  \item Olvidar la constante de integración $+C$.
\end{itemize}
\end{cuadroadvertencia}

\subsection{Triángulo de Pascal}

Se usa para expandir binomios de la forma $(a \pm b)^{n}$. Los coeficientes de la fila
$n$ se leen directamente; para $(a-b)^{n}$ los signos se \textbf{alternan}.

\[
\begin{array}{c@{\quad}l}
n=0 & 1\\[2pt]
n=1 & 1\;\;1\\[2pt]
n=2 & 1\;\;2\;\;1\\[2pt]
n=3 & 1\;\;3\;\;3\;\;1\\[2pt]
n=4 & 1\;\;4\;\;6\;\;4\;\;1\\[2pt]
n=5 & 1\;\;5\;\;10\;\;10\;\;5\;\;1\\[3pt]
n=6 & \boxed{1\;\;6\;\;15\;\;20\;\;15\;\;6\;\;1}\\[3pt]
n=7 & 1\;\;7\;\;21\;\;35\;\;35\;\;21\;\;7\;\;1
\end{array}
\]

\noindent Ejemplo (fila $n=6$):
\[
(u-2)^{6} = u^{6} - 12u^{5} + 60u^{4} - 160u^{3} + 240u^{2} - 192u + 64
\]

\subsection{División de polinomios}

\begin{cuadroteoria}
\textbf{Cuándo dividir.} Si el grado del numerador es \emph{mayor o igual} que el grado
del denominador, conviene dividir antes de integrar. Al dividir se obtiene un cociente
$Q(x)$ y un residuo $R(x)$ con grado menor que el del divisor:
\[
\frac{N(x)}{D(x)} = Q(x) + \frac{R(x)}{D(x)},
\qquad \deg R < \deg D
\]
\[
\Longrightarrow\qquad
\int \frac{N(x)}{D(x)}\dx = \int Q(x)\dx + \int \frac{R(x)}{D(x)}\dx
\]
El cociente $Q$ es un polinomio fácil de integrar término a término; la fracción
sobrante suele resolverse con un cambio de variable.
\end{cuadroteoria}

\begin{cuadroextra}
\textbf{Más adelante:} cuando el residuo $R(x)/D(x)$ no se resuelva por sustitución
directa, se estudiará el método de \textbf{fracciones parciales}. \emph{En esta guía
todavía no se usa}: todos los residuos se resuelven con cambio de variable.
\end{cuadroextra}

\newpage
\subsection{Repertorio trigonométrico (solo referencia)}

\begin{cuadroextra}
\small
Estas fórmulas aparecen en tus apuntes de clase rotuladas \texttt{m1}--\texttt{m6}
($a$ constante, $a \neq 0$). \textbf{Se mencionan como repertorio para más adelante};
en esta guía \emph{no} hay ejercicios con ellas.
\[
\begin{array}{ll}
\text{m1}:\ \displaystyle\int \sin(ax)\dx = -\frac{\cos(ax)}{a} \CC &
\text{m2}:\ \displaystyle\int \cos(ax)\dx = \frac{\sin(ax)}{a} \CC \\[10pt]
\text{m3}:\ \displaystyle\int \sec^{2}(ax)\dx = \frac{\tan(ax)}{a} \CC &
\text{m4}:\ \displaystyle\int \csc^{2}(ax)\dx = -\frac{\cot(ax)}{a} \CC \\[10pt]
\text{m5}:\ \displaystyle\int \sec(ax)\tan(ax)\dx = \frac{\sec(ax)}{a} \CC &
\text{m6}:\ \displaystyle\int \csc(ax)\cot(ax)\dx = -\frac{\csc(ax)}{a} \CC
\end{array}
\]
\end{cuadroextra}

\newpage
% ============================================================
\section{Bloque I — Cambio de variable directo}
% ============================================================

\begin{cuadroinfo}
\small
\textbf{Objetivo:} reconocer que un factor del integrando es (salvo constante) la
derivada del argumento interno. Resuelve con $u = $ \emph{el argumento interno}.
\end{cuadroinfo}

\vspace{4pt}

\begin{enumerate}[label=\textbf{\arabic*.}, leftmargin=*, itemsep=9pt]

  \item \medio\quad $\displaystyle\int x\,(x^{2}+5)^{6}\dx$

  \item \medio\quad $\displaystyle\int x^{2}\,(x^{3}-4)^{5}\dx$

  \item \medio\quad $\displaystyle\int (2x+1)\,(x^{2}+x+3)^{4}\dx$

  \item \medio\quad $\displaystyle\int \frac{x+1}{\sqrt{x^{2}+2x+5}}\dx$

  \item \medio\quad $\displaystyle\int \frac{x^{2}}{(x^{3}+7)^{2}}\dx$

  \item \medio\quad $\displaystyle\int (3x^{2}+2x)\,(x^{3}+x^{2}+1)^{10}\dx$

  \item \medio\quad $\displaystyle\int x^{3}\sqrt{x^{4}+2}\dx$

  \item \dificil\quad $\displaystyle\int \frac{x+2}{(x^{2}+4x+7)^{3}}\dx$

\end{enumerate}

\vspace{6pt}

% ============================================================
\section{Bloque II — Radicales y potencias fraccionarias}
% ============================================================

\begin{cuadroinfo}
\small
\textbf{Objetivo:} cuando la sustitución deja términos en $x$, despeja $x$ de
$u = g(x)$ y sustituye. Convierte raíces a exponentes fraccionarios y usa
$\int u^{k}\du = \dfrac{u^{k+1}}{k+1}+C$.
\end{cuadroinfo}

\vspace{4pt}

\begin{enumerate}[label=\textbf{\arabic*.}, leftmargin=*, itemsep=9pt, start=9]

  \item \dificil\quad $\displaystyle\int x\sqrt{x+3}\dx$

  \item \muyDificil\quad $\displaystyle\int x^{2}\sqrt{x-1}\dx$

  \item \dificil\quad $\displaystyle\int x\sqrt{2x+1}\dx$

  \item \dificil\quad $\displaystyle\int \frac{x}{\sqrt[3]{x+1}}\dx$

  \item \muyDificil\quad $\displaystyle\int x^{5}\sqrt{x^{3}+1}\dx$

  \item \dificil\quad $\displaystyle\int \frac{x^{3}}{\sqrt{x^{4}+1}}\dx$

  \item \dificil\quad $\displaystyle\int \frac{x^{2}}{\sqrt[3]{x^{3}+2}}\dx$

  \item \muyDificil\quad $\displaystyle\int \sqrt{1+\sqrt{x}}\dx$

  \item \muyDificil\quad $\displaystyle\int \frac{x^{3}}{(x^{2}+1)^{3/2}}\dx$

  \item \muyDificil\quad $\displaystyle\int \frac{x^{7}}{x^{4}+1}\dx$

\end{enumerate}

\vspace{6pt}

% ============================================================
\section{Bloque III — División de polinomios previa}
% ============================================================

\begin{cuadroinfo}
\small
\textbf{Objetivo:} si el grado del numerador $\geq$ el del denominador, divide y
escribe $\dfrac{N}{D} = Q + \dfrac{R}{D}$. Integra $Q$ término a término y resuelve
$\dfrac{R}{D}$ por sustitución.
\end{cuadroinfo}

\vspace{4pt}

\begin{enumerate}[label=\textbf{\arabic*.}, leftmargin=*, itemsep=9pt, start=19]

  \item \dificil\quad $\displaystyle\int \frac{x^{4}+x}{x+1}\dx$

  \item \dificil\quad $\displaystyle\int \frac{x^{5}-1}{x-1}\dx$

  \item \dificil\quad $\displaystyle\int \frac{x^{6}+2x^{5}+x^{4}+7}{(x+1)^{2}}\dx$

  \item \dificil\quad $\displaystyle\int \frac{4x^{4}+4x^{3}+x^{2}+6}{(2x+1)^{2}}\dx$

  \item \muyDificil\quad $\displaystyle\int \frac{x^{6}+2x^{4}+x^{2}+3x}{(x^{2}+1)^{2}}\dx$

  \item \muyDificil\quad $\displaystyle\int \frac{x^{7}+4x^{5}+4x^{3}+5x}{(x^{2}+2)^{2}}\dx$

  \item \muyDificil\quad $\displaystyle\int \frac{x^{6}-2x^{4}+x^{2}+4x}{(x^{2}-1)^{2}}\dx$

  \item \muyDificil\quad $\displaystyle\int \frac{x^{7}+6x^{5}+9x^{3}+4x}{(x^{2}+3)^{2}}\dx$

  \item \muyDificil\quad $\displaystyle\int \frac{x^{10}+2x^{7}+x^{4}+6x^{2}}{(x^{3}+1)^{2}}\dx$

  \item \muyDificil\quad $\displaystyle\int \frac{x^{7}+2x^{4}+3x^{2}+x}{(x^{3}+1)^{2}}\dx$

  \item \muyDificil\quad $\displaystyle\int \frac{x^{4}+4x^{2}+8x}{x^{2}+4}\dx$

  \item \muyDificil\quad $\displaystyle\int \frac{x^{8}+4x^{5}+13x^{2}}{(x^{3}+2)^{2}}\dx$

\end{enumerate}

\newpage
% ============================================================
\ifmwclaves
\section{Clave de Respuestas}
% ============================================================
Todas las respuestas incluyen la constante de integración $C \in \R$.

\begin{enumerate}[label=\textbf{\arabic*.}, leftmargin=*, itemsep=5pt]

  % ── Bloque I ─────────────────────────────────────────────
  \item $\dfrac{1}{14}(x^{2}+5)^{7} \CC$

  \item $\dfrac{1}{18}(x^{3}-4)^{6} \CC$

  \item $\dfrac{1}{5}(x^{2}+x+3)^{5} \CC$

  \item $\sqrt{x^{2}+2x+5} \CC$

  \item $-\dfrac{1}{3(x^{3}+7)} \CC$

  \item $\dfrac{1}{11}(x^{3}+x^{2}+1)^{11} \CC$

  \item $\dfrac{1}{6}(x^{4}+2)^{3/2} \CC$

  \item $-\dfrac{1}{4(x^{2}+4x+7)^{2}} \CC$

  % ── Bloque II ────────────────────────────────────────────
  \item $\dfrac{2}{5}(x+3)^{5/2} - 2(x+3)^{3/2} \CC$

  \item $\dfrac{2}{7}(x-1)^{7/2} + \dfrac{4}{5}(x-1)^{5/2} + \dfrac{2}{3}(x-1)^{3/2} \CC$

  \item $\dfrac{1}{10}(2x+1)^{5/2} - \dfrac{1}{6}(2x+1)^{3/2} \CC$

  \item $\dfrac{3}{5}(x+1)^{5/3} - \dfrac{3}{2}(x+1)^{2/3} \CC$

  \item $\dfrac{2}{15}(x^{3}+1)^{5/2} - \dfrac{2}{9}(x^{3}+1)^{3/2} \CC$

  \item $\dfrac{1}{2}\sqrt{x^{4}+1} \CC$

  \item $\dfrac{1}{2}(x^{3}+2)^{2/3} \CC$

  \item $\dfrac{4}{5}(1+\sqrt{x})^{5/2} - \dfrac{4}{3}(1+\sqrt{x})^{3/2} \CC$

  \item $\dfrac{x^{2}+2}{\sqrt{x^{2}+1}} \CC$

  \item $\dfrac{x^{4}}{4} - \dfrac{1}{4}\ln(x^{4}+1) \CC$

  % ── Bloque III ───────────────────────────────────────────
  \item $\dfrac{x^{4}}{4} - \dfrac{x^{3}}{3} + \dfrac{x^{2}}{2} \CC$

  \item $\dfrac{x^{5}}{5} + \dfrac{x^{4}}{4} + \dfrac{x^{3}}{3} + \dfrac{x^{2}}{2} + x \CC$

  \item $\dfrac{x^{5}}{5} - \dfrac{7}{x+1} \CC$

  \item $\dfrac{x^{3}}{3} - \dfrac{3}{2x+1} \CC$

  \item $\dfrac{x^{3}}{3} - \dfrac{3}{2(x^{2}+1)} \CC$

  \item $\dfrac{x^{4}}{4} - \dfrac{5}{2(x^{2}+2)} \CC$

  \item $\dfrac{x^{3}}{3} - \dfrac{2}{x^{2}-1} \CC$

  \item $\dfrac{x^{4}}{4} - \dfrac{2}{x^{2}+3} \CC$

  \item $\dfrac{x^{5}}{5} - \dfrac{2}{x^{3}+1} \CC$

  \item $\dfrac{x^{2}}{2} - \dfrac{1}{x^{3}+1} \CC$

  \item $\dfrac{x^{3}}{3} + 4\ln(x^{2}+4) \CC$

  \item $\dfrac{x^{3}}{3} - \dfrac{3}{x^{3}+2} \CC$

\end{enumerate}
\fi

\end{document}
